\section{Conclusiones de ingeniería, limitaciones y preguntas del curso}

El último grupo de diapositivas condensa los resultados de la investigación en principios de ingeniería: recompensar los pensamientos intermedios, usar señales densas y verificar la eficiencia de los tokens, combinando además el currículo, el razonamiento front-loading y el reforzamiento. Sin embargo, las conclusiones realmente aplicables requieren pasar por auditoría de datos, contabilidad de costos, análisis de fallos y límites de seguridad.

\subsection{Tres ablationes: por qué funciona, cuán densas deben ser las señales y cuántos datos ahorrar}

La página de Key Ablations desglosa el valor del RLP en tres puntos: Reward to Think verifica si los pensamientos obtienen crédito directo; Sparse vs Dense examina la diferencia entre recompensar solo el resultado correcto final y la ganancia de información por posición; Token Efficiency evalúa la utilización de datos relativa a una línea base NTP enorme. Las tres columnas demuestran conjuntamente que los beneficios provienen del diseño del objetivo, no simplemente de generar más segmentos de CoT.

\lecturefigure{slide-041.jpg}{Conclusiones de reward-to-think, recompensa densa y eficiencia de tokens del RLP}{Grabación oficial de Stanford, 00:45:56--00:46:27}

\begin{warningbox}{Las ablationes aún requieren una tabla de costos completa}
Si la recompensa densa utiliza más rollouts, más forwards de scorer o pensamientos más largos, también altera la señal y el cómputo. Una ablación ideal debería reportar los mismos FLOPs, el mismo tiempo wall-clock, el mismo número de tokens o varios métodos de emparejamiento, evitando confundir el cómputo adicional con el efecto real de la forma de recompensa.
\end{warningbox}

\subsection{De Pascal a Hopper: las capacidades provienen de la superposición de estrategias}

La curva reconecta los cuatro learners: el currículo aporta la primera capa de mejora, el front-loading del razonamiento eleva más el fundamento y, al agregar RLP con learning through thinking, se alcanza la mayor mejora relativa. La curva es una síntesis pedagógica a través de experimentos, no una descomposición causal estrictamente aditiva en un único entrenamiento unificado; por lo tanto, no se deben tomar los porcentajes de cada tramo como constantes independientes y superponibles.

\lecturefigure{slide-042.jpg}{De Pascal a Hopper: agregar gradualmente tres estrategias de aprendizaje}{Grabación oficial de Stanford, 00:46:27--00:47:12}

\begin{importantbox}{Juicio integral en lugar de suma mecánica}
El currículo cambia la distribución que ve el modelo, el front-loading altera la representación base y RLP modifica nuevamente el paisaje de optimización de la política de pensamiento; existen interacciones entre ellos. En una verdadera combinación se debe realizar una ablación factorial, no sumar los headline gains de cada artículo para predecir el beneficio final.
\end{importantbox}

\subsection{Seis conclusiones y tres no-conclusiones}

El curso concluye finalmente: dos fases son efectivas; el front-loading genera una ventaja duradera y compuesta; las señales de alta calidad pueden manifestarse después del SFT; RLP anticipa el razonamiento por reforzamiento; los textos sin etiquetar estándar también pueden proporcionar señales de entrenamiento tipo razonamiento; y la exploración junto con la recompensa abren un nuevo eje para el escalado del razonamiento. Cada una debe ir acompañada de límites de modelo, datos y evaluación.

\lecturefigure{slide-043.jpg}{Currículo, razonamiento y reforzamiento redefinen el preentrenamiento}{Grabación oficial de Stanford, 00:47:12--00:49:00}

En la figura \textbf{no se demuestra} nada de esto: en primer lugar, que todas las cadenas de pensamiento merecen ser retenidas; en segundo, que RLP reduciría automáticamente las alucinaciones o mejoraría el alineamiento; y en tercer lugar, que los resultados solo de lenguaje transferirían directamente a un modelo de visión-lenguaje. El orador aclara explícitamente en Q\&A que multimodal se considera no testeado, mientras que RLHF/alucinaciones son temas adyacentes pero no parte de este trabajo.

\subsection{La receta final vuelve a la colaboración sistémica}

La página de cierre muestra nuevamente los cuatro componentes, pero ahora cada columna tiene un significado concreto: Smart Data requiere etiquetas de calidad/dominio rastreables; Smart Architecture necesita soporte para entrenamiento largo y rollouts; Smart Algorithms incluye currículo, front-loading y RLP; Smart Collaboration se encarga de unificar las métricas de preentrenamiento/post-entrenamiento y el presupuesto de costos. Recuerda al lector que los artículos de algoritmos solo pueden convertirse en capacidades del modelo al ingresar a un pipeline completo.

\lecturefigure{slide-044.jpg}{Revisar nuevamente la receta de los cuatro componentes tras aprender mecanismos de aprendizaje}{Grabación oficial de Stanford, 00:49:00--00:49:30}

\begin{knowledgebox}{Un orden de decisiones ejecutable}
Primero audite los datos y las métricas de evaluación, luego use un proxy a pequeña escala para estimar calidad/repetición/límite de fase; posteriormente pruebe el front-loading con un presupuesto fijo; solo cuando la línea base, los FLOPs y el flujo post-entrenamiento estén estables, introduzca los rollouts de RLP. En cada paso conserve las opciones no adoptadas y los resultados fallidos, evitando convertir el ajuste de parámetros de la receta en alquimia irreproducible.
\end{knowledgebox}

\subsection{Preguntas del curso: detalles de implementación y límites}

Las preguntas complementan varios puntos clave. La maquinaria de ventaja de RLP es similar a GRPO; la principal variación es la recompensa de ganancia de información. Los experimentos seleccionan tokens aleatoriamente en el documento, no basándose en un cribado preciso por entropía. RLP usa una mezcla general de preentrenamiento, incluyendo gran cantidad de web crawl, y no opera solo sobre datos de razonamiento de alta calidad. También se observaron beneficios relativos a CPT desde checkpoints iniciales del orden del 20\% del preentrenamiento completo.

\teachervoice{00:49:54--00:50:35, el orador explica que RLP reutiliza la ventaja relativa grupal estilo GRPO; el foco de innovación es la construcción de recompensa.}

\teachervoice{00:55:36--00:57:43, el orador corrige al preguntante: RLP no funciona solo sobre datos de más alta calidad; en la práctica se seleccionan tokens aleatoriamente y también se realizaron experimentos favorables en checkpoints tempranos del orden de 4 T-token.}

\begin{lstlisting}[caption={Tabla mínima de experimento recomendada para evaluar RLP}]
experiment = {
    "start_checkpoint_tokens": ...,   # e.g. 19.8T
    "extra_data_tokens": ...,
    "rollouts_per_position": ...,
    "mean_thought_length": ...,
    "estimated_training_flops": ...,
    "gpu_hours": ...,
    "token_selection": "uniform_random",
    "data_mixture_version": ...,
    "post_training_recipe": ...,
    "benchmark_harness": ...,
}
\end{lstlisting}

\subsection{Límites de calidad de datos, RLHF y multimodal}

Sobre la calidad general de las páginas web, el orador menciona un clasificador educativo estilo FineWeb-Edu: usar una rúbrica para dar puntajes de 1--5 a los documentos, escalable además al dominio y nivel educativo. Respecto al sesgo de RLHF, alucinaciones y hacking de recompensa, ella indica claramente que es investigación continua de alineamiento, no una conclusión directa de las tres trabajos de esta conferencia. Para el alineamiento visión-lenguaje, su respuesta es que aún no se ha testeado.

\begin{warningbox}{No reescriba las respuestas cautelosas del Q\&A como conclusiones definitivas}
"RLHF aún puede ocupar una pequeña parte del post-entrenamiento", "el clasificador necesita actualización continua", "RLAIF podría reducir cierta subjetividad" son discusiones de clase y juicios empíricos, no validaciones experimentales de esta conferencia. Las apuntes mantienen sus límites, evitando escribir opiniones de áreas adyacentes como declaraciones de efectos de RLP.
\end{warningbox}

\subsection{Resumen de la sección}

Lo más importante desde el punto de vista de ingeniería no es repetir los headline gains, sino establecer un plano de control completo: calidad de datos y currículo, asignación de razonamiento, recompensa de rollout, estabilidad EMA, coincidencia token/FLOP/wall-clock, persistencia post-entrenamiento y restricciones transdominio. Solo con toda esta información disponible se puede juzgar si RLP o front-loading valen la pena ingresar al entrenamiento de producción.
